\section{Threats to Validity and Ethics}
\label{sec:threats}

\subsection{Labelling}

Both corpora label by a proxy. 5G Traffic inherits its label from the capture
directory, so background OS traffic, DNS and shared CDN fetches recorded during
a session are labelled as the foreground application; we do not filter them,
because a filter would need ground truth we lack. CESNET-QUIC22 derives its
labels from SNI, which makes the task \emph{predict what SNI would have said}
and caps the achievable ceiling at how well SNI itself categorises traffic. We
measure that ceiling rather than assuming it: 0.967 macro-F1 from the SNI
string alone.

Encrypted Client Hello~\cite{rescorla2024ech} removes this labelling mechanism.
That strengthens the motivation for traffic analysis without name resolution
and weakens future data collection in the same movement; corpora of this kind
may not be reproducible in kind.

\subsection{Residual dependence}

Our protocols control one axis each and leave others uncontrolled. On 5G
Traffic, segments from different captures of the same application share a
device and access network. On CESNET, we group by capture day but not by
client; a client-disjoint protocol is possible, since \texttt{SRC\_IP} is
retained, and is the obvious next control. We report what each protocol does
and does not separate rather than implying a partition is clean.

As noted in \S\ref{sec:datasets}, 5G session-disjoint splitting is partly
application-disjoint by accident, because several applications have a single
capture file. Six of fifteen applications have no test rows under that split
and one has no training rows. We exclude zero-support classes from macro
averages and report how many were excluded.

\subsection{Sampling and budget}

CESNET rows are a seeded reservoir of 400 flows per class per day over all 28
days: every row is considered, but 201{,}600 of roughly 100 million are
retained, and class balance is imposed at read time. Reported class priors are
therefore ours, not the network's, and no claim about base rates follows from
this work. Ablations run at a reduced budget shared by every row of their
family, and the ablation table declares itself partial: the component and
hyperparameter sweeps were not run, because the closed-set result they would
explain does not hold.

\subsection{Statistical power}

Our largest effect --- server-disjoint versus session-disjoint, Cohen's
$d = 20.3$ --- reaches $p = 0.031$ at six seeds and does not survive
Holm-Bonferroni correction across its three-comparison family, which requires
$p \le 0.017$. This is arithmetic, not weakness: a two-sided Wilcoxon
signed-rank test on $n$ paired seeds has a minimum achievable $p$-value of
$2^{-(n-1)}$, so six seeds cannot certify \emph{any} effect in a family of
three. Eight seeds are required and were run. We note that a results section
reporting corrected Wilcoxon tests over five seeds --- a common protocol ---
cannot report a corrected significant result at all.

\subsection{Generalisation}

Two corpora, two networks, one protocol family. Fifteen applications on one and
roughly a hundred on the other; neither approaches the open world. Both are
from 2022.

\subsection{Ethics}

\paragraph{Dual use.} This work concerns identifying applications from
encrypted traffic metadata. The capability serves network operations and it
serves censorship and surveillance. We take that seriously rather than treating
it as incidental. Our principal findings cut \emph{against} an operational use:
we show these classifiers are far less reliable than the literature reports
once evaluation controls for how the data was collected, and we quantify the
padding overhead that defeats them. Overstating their reliability would be the
more harmful error.

\paragraph{Data.} We collected nothing. Both corpora are public and published
for research under CC BY 4.0. CESNET-QUIC22 is anonymised at source by its
publishers under an institutional data-handling policy. The 5G captures are of
the original authors' own devices. We retain client addresses solely as
split-grouping keys, never as model inputs, and a continuous-integration test
asserts that mechanically. No raw capture is redistributed; our released
artifacts are split manifests, metric files and code.

\paragraph{Defender-facing disclosure.} Our countermeasure evaluation reports
accuracy degradation against the overhead each defence imposes, so that the
cost of defeating this class of classifier is stated publicly. We find that the
cheap defences --- random padding, dummy injection, size quantisation --- cost
a defender 8--17\% overhead and cost an attacker essentially nothing, while
only defences costing 2.28$\times$ bandwidth degrade classification
meaningfully. A defender deploying the cheap options is paying for nothing, and
that is information defenders need.
